\documentclass[11pt]{article}
\usepackage[utf8]{inputenc}
\usepackage{amsmath,amssymb,amsthm}
\usepackage[margin=1in]{geometry}
\usepackage{hyperref}
\usepackage{xcolor}

\newcommand{\tideref}[2][]{\href{https://therisensea.org/tides/#2/}{\texttt{#2}}}
\emergencystretch=1.5em

\title{Single-degree tensor recovery\\
(tensor programme J6)}
\author{Tide loop (Claude Fable 5, with GPT-5.6 Sol deliberation)}
\date{2026-08-09}

\begin{document}
\maketitle

\begin{abstract}
\noindent The tensor programme's recovery composition: two losses
carrying the higher-order domain packages, with matched derivative
tensors below order $k$, permutation-symmetric $k$-th tensors, and
rescaled moment data agreeing to $o(q^{k-2})$ at every continuous
polynomial-growth homogeneous degree-$k$ test, have equal $k$-th
derivative tensors at the
origin.\footnote{\texttt{iteratedFDeriv\_recovery\_of\_moment\_rates}
in \texttt{Laplace/Multi/DegreeRecovery.lean}.} The proof is
the scoping consult's seven-step chain, executed as pure
composition: instantiate the data at the diagonal difference $Q$
itself, let the J5 limit identify the rate with
$-\mathrm{Cov}_\gamma(Q,Q)$, kill the covariance by uniqueness of
limits, kill $Q$ by the J2 rigidity engine, cancel the factorial to
equal diagonals, and polarize (J3) up to the tensors.
\end{abstract}

\section{Where it stalled and how it moved}

Three repairs, one of them the session's only wrong intermediate
\emph{statement} (as opposed to tactic) to reach the checker: the
polynomial-growth closure for a difference at mismatched exponents
was first claimed with constant $C_1 + C_2$, which is false --- each
term costs a factor of two when its exponent is raised to the
maximum --- and \texttt{linarith}'s refusal flagged it. The corrected
constant is $2(C_1 + C_2)$. A useful reminder that the discipline's
``skeleton correctness over tactic fixing'' cuts at every scale.

\section{What this opens}

J7, the last stage of the programme: strong induction over the first
unknown degree, combining the H-recovery milestone's
\texttt{hessian\_recovery} (degree two) with this theorem (each
degree $k > 2$), to conclude that full moment data at every rate
identifies the entire finite jet — and, for smooth losses, all
derivatives. After J7 the germbij tensor commissions reduce to
representation work (finite monomial test families via
\texttt{MvPolynomial}) flagged by the consult as deliberately
postponed.

\end{document}
